% Zone „Kennst du schon" – Kurvenuntersuchung, Kl. 12 Berlin, Gymnasium GK
\setcounter{aufgabe}{0}
\einheitenkopf[zone][\zonekurz]{Das kennst du schon}

\begin{aufgabe}{Ich kann Hochpunkte, Tiefpunkte und Achsenschnittpunkte an einem Graphen ablesen.}
Lies am Graphen von $f$ ab.

\begin{ksys}[xmin=-5,xmax=4,ymin=-2,ymax=5,ablesen]
  \funktion{0.125*\x^3-1.5*\x+2}{f}
\end{ksys}

\begin{teile}
\teil Schnittpunkt mit der $y$-Achse: \leerfeld
\teil Hochpunkt: \leerfeld
\teil Tiefpunkt: \leerfeld
\teil Schnittpunkte mit der $x$-Achse: \leerfeld
\teil An welcher Stelle berührt der Graph die $x$-Achse, ohne sie zu schneiden? \leerfeld
\end{teile}
\end{aufgabe}

\begin{aufgabe}{Ich kann Gleichungen durch Ausklammern und mit der Lösungsformel lösen.}
Löse die Gleichung.
\begin{gleichungsraster}[2]
\gl{4x - 12 = 0} & \gl{x\,(x - 5) = 0} \\
\gl{x^2 - 6x + 5 = 0} & \gl{x^2 - 36 = 0} \\
\gl{6x^2 + 6x - 36 = 0} & \gl{x^3 - 4x = 0} \\
\gl{x^4 - 5x^2 = 0} & \gl{(x - 2)\cdot e^{x} = 0} \\
\end{gleichungsraster}
\end{aufgabe}

\begin{aufgabe}{Ich finde den Fehler beim Lösen einer Gleichung durch Teilen.}
Lena löst die Gleichung $x^3 = 16x$ so:
\rechnung{x^3 &= 16x &&\mid :x \\ x^2 &= 16 \\ x &= -4 \text{ oder } x = 4}
\begin{teile}
\teil Finde den Fehler, benenne ihn und korrigiere die Rechnung.
\end{teile}
\end{aufgabe}

\begin{aufgabe}{Ich kann eine Gleichung lösen, ohne durch $x$ zu teilen.}
Löse die Gleichung.
\begin{gleichungsraster}[3]
\gl{2x^3 = 18x} & \gl{x^4 = 7x^3} \\
\end{gleichungsraster}
\end{aufgabe}

\begin{aufgabe}{Ich kann Ableitungen bilden.}
Bilde die erste Ableitung.
\begin{gleichungsraster}[1]
\gl{f(x) = x^4} & \gl{f(x) = 5x^2} \\
\gl{f(x) = 2x^3 - 6x^2 + 7} & \gl{f(x) = -\tfrac{1}{4}x^4 + 2x} \\
\gl[2]{f(x) = e^{2x}} & \gl[2]{f(x) = x\cdot e^{x}} \\
\gl[2]{f(x) = (x - 3)\cdot e^{-x}} & \\
\anweisung{Bilde die erste und die zweite Ableitung.}
\gl[2]{f(x) = x^4 - 2x^3} & \gl[3]{f(x) = x\cdot e^{2x}} \\
\end{gleichungsraster}
\end{aufgabe}

\begin{aufgabe}{Ich kann Funktionswerte berechnen und die Punktprobe machen.}
Gegeben sind $f(x) = x^3 - 2x^2 + 1$ und $g(x) = (x + 1)\cdot e^{x}$. Berechne den Funktionswert.
\begin{teilezwei}
\tz $f(0)$ = \leerfeld & \tz $f(2)$ = \leerfeld \\
\tz $f(-1)$ = \leerfeld & \tz $f(0{,}5)$ = \leerfeld \\
\tz $g(0)$ = \leerfeld & \tz $g(-1)$ = \leerfeld \\
\anweisung{Prüfe, ob der Punkt auf dem Graphen von $f$ liegt.}
\tz $P(2\,|\,1)$ \quad \janein & \tz $Q(-2\,|\,9)$ \quad \janein \\
\end{teilezwei}
\end{aufgabe}

\begin{aufgabe}{Ich kann den Verlauf eines Graphen am Funktionsterm grob skizzieren.}
Skizziere den Graphen grob in das Koordinatensystem mit demselben Buchstaben. Trage zuerst die Nullstellen ein.
\begin{teile}
\teil $f(x) = (x - 2)(x + 1)$
\teil $g(x) = -x^2 + 4$
\teil $h(x) = -x^3 + 4x$
\end{teile}

\begin{ksys}[xmin=-3,xmax=3,ymin=-4,ymax=4]
  \node[anchor=north west,font=\bfseries] at (-3,4) {a)};
\end{ksys}\ksysabstand\begin{ksys}[xmin=-3,xmax=3,ymin=-4,ymax=4]
  \node[anchor=north west,font=\bfseries] at (-3,4) {b)};
\end{ksys}\ksysabstand\begin{ksys}[xmin=-3,xmax=3,ymin=-4,ymax=4]
  \node[anchor=north west,font=\bfseries] at (-3,4) {c)};
\end{ksys}

\end{aufgabe}

\begin{aufgabe}{Ich kann die Ableitung als Steigung des Graphen und als Änderungsrate deuten.}
Gib an, ob der Graph von $q$ im Punkt steigt, fällt oder waagerecht verläuft und ob $q'$ dort positiv, negativ oder null ist.

\begin{ksys}[xmin=-1,xmax=6,ymin=-2,ymax=4,ablesen]
  \funktion{0.5*\x^2-2*\x+1}{q}
  \tangentean*{0.5*\x^2-2*\x+1}{0}{s}
  \punkt{0}{1}{A}
  \punkt{2}{-1}{B}
  \punkt{4}{1}{C}
\end{ksys}

\begin{teile}
\teil im Punkt $A$: \leerfeld
\teil im Punkt $B$: \leerfeld
\teil im Punkt $C$: \leerfeld
\anweisung{Lies ab und entscheide.}
\teil Steigung der Tangente $s$ in $A$: \leerfeld
\teil $A$ und $C$ haben denselben Funktionswert. Haben $q'(0)$ und $q'(4)$ auch denselben Wert? Begründe.
\anweisung{Deute im Sachzusammenhang.}
\teil $V(t)$ ist die Wassermenge in einem Becken in Litern, $t$ die Zeit in Minuten. Es gilt $V'(10) = -4$. Was bedeutet das?
\end{teile}
\end{aufgabe}
